% GENERATED FILE. Edit canonical lessons, then run npm run generate:books.
% canonical-source-hash: fnv1a64:ee1ec2e8cbc67f6a
% canonical-lessons: GU-C385-puraskar, GU-C385-kadar, GU-C385-mahentanu, GU-C385-pamvu, GU-C385-navajvu, GU-R385-review-379-385

\chapter{Rewards}
\label{ch:gu-rewards}
\hlchaptermodality{385}

\begin{chapteropening}
\textbf{By the end of this chapter:} \emph{I can say what someone will get if they do well.}

Complete the last lesson of chapter 385: I can say what someone will get if they do well.
\end{chapteropening}

\section[\texorpdfstring{\gu{પુરસ્કાર} (puraskār)}{પુરસ્કાર (puraskār)}]{\gu{પુરસ્કાર} (puraskār) — an award}

\label{lesson:GU-C385-puraskar}



\begin{quote}
Before the new one: say the Gujarati for to get one's way, then the Gujarati for to deliver.
\end{quote}

\subsection*{You'll want to know: \gu{પુરસ્કાર}}
\textbf{\gu{પુરસ્કાર}} — \emph{puraskār} — \textquotedblleft{}an award\textquotedblright{}.

\begin{quote}
\gu{જો} \gu{મીરા} \gu{હરીફાઈ} \gu{જીતશે} \gu{તો} \gu{શાળા} \gu{તેને} \gu{પુરસ્કાર} \gu{આપશે}.
\end{quote}

\begin{quote}
\emph{jo Mīrā harīphāī jītśe to śāḷā tene puraskār āpśe.}
\end{quote}

\begin{quote}
If Mira wins the contest, the school will give her an award.
\end{quote}

A masculine noun.

\begin{cousinweb}
\textbf{\gu{પુરસ્કાર}} (\emph{puraskār}) is a tatsama, borrowed from Sanskrit \emph{puraskāra}, \textquotedblleft{}placing in front, honouring\textquotedblright{} (Monier-Williams).
\end{cousinweb}

\begin{grammarlens}[title={What you've built}]
One more word for saying what happens if.
\end{grammarlens}

\subsection*{Your turn}
\begin{itemize}
\raggedright
  \item \emph{Say it:} \emph{puraskār}
  \item \emph{Say it:} \emph{puraskār}, once more
  \item \emph{Say it:} \emph{jo Mīrā harīphāī jītśe to śāḷā tene puraskār āpśe.} — If Mira wins the contest, the school will give her an award.
  \item \emph{Recall:} say the Gujarati for to get one's way, then the Gujarati for to deliver, then say \emph{puraskār} again
\end{itemize}

\begin{tcolorbox}[breakable,colback=teal!4,colframe=teal!35!black,title={Before you move on}]
What does \gu{પુરસ્કાર} mean? (\textquotedblleft{}An award\textquotedblright{}.) Say it once more.
\end{tcolorbox}
\section[\texorpdfstring{\gu{કદર} (kadar)}{કદર (kadar)}]{\gu{કદર} (kadar) — appreciation}

\label{lesson:GU-C385-kadar}



\begin{quote}
Before the new one: say the Gujarati for to deliver, then the Gujarati for an award.
\end{quote}

\subsection*{You'll want to know: \gu{કદર}}
\textbf{\gu{કદર}} — \emph{kadar} — \textquotedblleft{}appreciation\textquotedblright{}.

\begin{quote}
\gu{જો} \gu{સાહેબ} \gu{તમારા} \gu{કામની} \gu{કદર} \gu{કરશે} \gu{તો} \gu{તમારો} \gu{પગાર} \gu{વધશે}.
\end{quote}

\begin{quote}
\emph{jo sāheb tamārā kāmnī kadar karśe to tamāro pagār vadhśe.}
\end{quote}

\begin{quote}
If the boss appreciates your work, your pay will go up.
\end{quote}

A feminine noun.

\begin{cousinweb}
\textbf{\gu{કદર}} (\emph{kadar}) is usually traced to Arabic \emph{qadr}, \textquotedblleft{}worth, value\textquotedblright{}, through Persian.
\end{cousinweb}

\begin{grammarlens}[title={What you've built}]
One more word for saying what happens if.
\end{grammarlens}

\subsection*{Your turn}
\begin{itemize}
\raggedright
  \item \emph{Say it:} \emph{kadar}
  \item \emph{Say it:} \emph{kadar}, once more
  \item \emph{Say it:} \emph{jo sāheb tamārā kāmnī kadar karśe to tamāro pagār vadhśe.} — If the boss appreciates your work, your pay will go up.
  \item \emph{Recall:} say the Gujarati for to deliver, then the Gujarati for an award, then say \emph{kadar} again
\end{itemize}

\begin{tcolorbox}[breakable,colback=teal!4,colframe=teal!35!black,title={Before you move on}]
What does \gu{કદર} mean? (\textquotedblleft{}Appreciation\textquotedblright{}.) Say it once more.
\end{tcolorbox}
\section[\texorpdfstring{\gu{મહેનતાણું} (mahentāṇũ)}{મહેનતાણું (mahentāṇũ)}]{\gu{મહેનતાણું} (mahentāṇũ) — a fee (payment for work)}

\label{lesson:GU-C385-mahentanu}



\begin{quote}
Before the new one: say the Gujarati for an award, then the Gujarati for appreciation.
\end{quote}

\subsection*{You'll want to know: \gu{મહેનતાણું}}
\textbf{\gu{મહેનતાણું}} — \emph{mahentāṇũ} — \textquotedblleft{}a fee (payment for work)\textquotedblright{}.

\begin{quote}
\gu{જો} \gu{કામ} \gu{સારું} \gu{થશે} \gu{તો} \gu{હું} \gu{વધારે} \gu{મહેનતાણું} \gu{ચૂકવીશ}.
\end{quote}

\begin{quote}
\emph{jo kām sārũ thaśe to hũ vadhāre mahentāṇũ cūkvīś.}
\end{quote}

\begin{quote}
If the work is good, I will pay a higher fee.
\end{quote}

A neuter noun.

\begin{cousinweb}
\textbf{\gu{મહેનતાણું}} (\emph{mahentāṇũ}) is built on \emph{mahenat}, \textquotedblleft{}hard work\textquotedblright{}, usually traced to Arabic \emph{miḥnat}, \textquotedblleft{}trial, toil\textquotedblright{}, through Persian, with the ending \emph{-āṇũ}: what labour earns.
\end{cousinweb}

\begin{grammarlens}[title={What you've built}]
One more word for saying what happens if.
\end{grammarlens}

\subsection*{Your turn}
\begin{itemize}
\raggedright
  \item \emph{Say it:} \emph{mahentāṇũ}
  \item \emph{Say it:} \emph{mahentāṇũ}, once more
  \item \emph{Say it:} \emph{jo kām sārũ thaśe to hũ vadhāre mahentāṇũ cūkvīś.} — If the work is good, I will pay a higher fee.
  \item \emph{Recall:} say the Gujarati for an award, then the Gujarati for appreciation, then say \emph{mahentāṇũ} again
\end{itemize}

\begin{tcolorbox}[breakable,colback=teal!4,colframe=teal!35!black,title={Before you move on}]
What does \gu{મહેનતાણું} mean? (\textquotedblleft{}A fee (payment for work)\textquotedblright{}.) Say it once more.
\end{tcolorbox}
\section[\texorpdfstring{\gu{પામવું} (pāmvũ)}{પામવું (pāmvũ)}]{\gu{પામવું} (pāmvũ) — to attain}

\label{lesson:GU-C385-pamvu}



\begin{quote}
Before the new one: say the Gujarati for appreciation, then the Gujarati for payment for work.
\end{quote}

\subsection*{You'll want to know: \gu{પામવું}}
\textbf{\gu{પામવું}} — \emph{pāmvũ} — \textquotedblleft{}to attain\textquotedblright{}.

\begin{quote}
\gu{જો} \gu{તે} \gu{સતત} \gu{મહેનત} \gu{કરશે} \gu{તો} \gu{બધાનો} \gu{આદર} \gu{પામશે}.
\end{quote}

\begin{quote}
\emph{jo te satat mahenat karśe to badhāno ādar pāmśe.}
\end{quote}

\begin{quote}
If he keeps working hard, he will win everyone's respect.
\end{quote}

Past (he …): \gu{તે} \gu{બધાનો} \gu{આદર} \gu{પામ્યો}, \emph{te badhāno ādar pāmyo}.

\begin{cousinweb}
\textbf{\gu{પામવું}} (\emph{pāmvũ}) is inherited from Sanskrit \emph{prāpnōti}, \textquotedblleft{}attains to, obtains\textquotedblright{}; CDIAL lists the Gujarati verb (CDIAL 8947).
\end{cousinweb}

\begin{grammarlens}[title={What you've built}]
One more word for saying what happens if.
\end{grammarlens}

\subsection*{Your turn}
\begin{itemize}
\raggedright
  \item \emph{Say it:} \emph{pāmvũ}
  \item \emph{Say it:} \emph{pāmvũ}, once more
  \item \emph{Say it:} \emph{jo te satat mahenat karśe to badhāno ādar pāmśe.} — If he keeps working hard, he will win everyone's respect.
  \item \emph{Recall:} say the Gujarati for appreciation, then the Gujarati for payment for work, then say \emph{pāmvũ} again
\end{itemize}

\begin{tcolorbox}[breakable,colback=teal!4,colframe=teal!35!black,title={Before you move on}]
What does \gu{પામવું} mean? (\textquotedblleft{}To attain\textquotedblright{}.) Say it once more.
\end{tcolorbox}
\section[\texorpdfstring{\gu{નવાજવું} (navājvũ)}{નવાજવું (navājvũ)}]{\gu{નવાજવું} (navājvũ) — to honour with an award}

\label{lesson:GU-C385-navajvu}



\begin{quote}
Before the new one: say the Gujarati for payment for work, then the Gujarati for to attain.
\end{quote}

\subsection*{You'll want to know: \gu{નવાજવું}}
\textbf{\gu{નવાજવું}} — \emph{navājvũ} — \textquotedblleft{}to honour with an award\textquotedblright{}.

\begin{quote}
\gu{જો} \gu{તે} \gu{જીતશે} \gu{તો} \gu{સરકાર} \gu{તેને} \gu{પુરસ્કારથી} \gu{નવાજશે}.
\end{quote}

\begin{quote}
\emph{jo te jītśe to sarkār tene puraskārthī navājśe.}
\end{quote}

\begin{quote}
If he wins, the government will honour him with an award.
\end{quote}

Past, with \gu{તેણે} (he or she …): \gu{તેણે} \gu{ખેલાડીને} \gu{પુરસ્કારથી} \gu{નવાજ્યો}, \emph{teṇe khelāḍīne puraskārthī navājyo}.

\begin{cousinweb}
\textbf{\gu{નવાજવું}} (\emph{navājvũ}) is usually traced to Persian \emph{navāz}, \textquotedblleft{}caressing, favouring\textquotedblright{}, with the Gujarati verb ending.
\end{cousinweb}

\begin{grammarlens}[title={What you've built}]
One more word for saying what happens if.
\end{grammarlens}

\subsection*{Your turn}
\begin{itemize}
\raggedright
  \item \emph{Say it:} \emph{navājvũ}
  \item \emph{Say it:} \emph{navājvũ}, once more
  \item \emph{Say it:} \emph{jo te jītśe to sarkār tene puraskārthī navājśe.} — If he wins, the government will honour him with an award.
  \item \emph{Recall:} say the Gujarati for payment for work, then the Gujarati for to attain, then say \emph{navājvũ} again
\end{itemize}

\begin{tcolorbox}[breakable,colback=teal!4,colframe=teal!35!black,title={Before you move on}]
What does \gu{નવાજવું} mean? (\textquotedblleft{}To honour with an award\textquotedblright{}.) Say it once more.
\end{tcolorbox}
\section[(dialogue)]{Looking back: from Very, Extremely to Rewards}

\label{lesson:GU-R385-review-379-385}



\begin{quote}
Say the Gujarati for extremely and utterly.
\end{quote}

\subsection*{Your turn}
Say these aloud:

\begin{itemize}
\raggedright
  \item extremely, utterly, complete, absolutely, solely
  \item enough, a shortage, a little, slightly, to cope
  \item to regret, a pity, if only, to long for someone, an opening
  \item a goal, a resolution, an achievement, planning, to strive
  \item an attempt, an effort, enthusiasm, to set about, to battle on
  \item successful, advancement, development, to get one's way, to deliver
  \item an award, appreciation, payment for work, to attain, to honour with an award
\end{itemize}

\begin{tcolorbox}[breakable,colback=teal!4,colframe=teal!35!black,title={Before you move on}]
Extremely? (\textbf{\gu{અત્યંત}}, \emph{atyant}.) Utterly? (\textbf{\gu{સાવ}}, \emph{sāv}.) Complete? (\textbf{\gu{સંપૂર્ણ}}, \emph{sampūrṇ}.) Absolutely? (\textbf{\gu{બિલકુલ}}, \emph{bilkul}.) Solely? (\textbf{\gu{કેવળ}}, \emph{kevaḷ}.) Enough? (\textbf{\gu{પૂરતું}}, \emph{pūrtũ}.) A shortage? (\textbf{\gu{તંગી}}, \emph{tangī}.) A little? (\textbf{\gu{થોડું}}, \emph{thoḍũ}.) Slightly? (\textbf{\gu{સહેજ}}, \emph{sahej}.) To cope? (\textbf{\gu{પહોંચી} \gu{વળવું}}, \emph{pahõcī vaḷvũ}.) To regret? (\textbf{\gu{પસ્તાવું}}, \emph{pastāvũ}.) A pity? (\textbf{\gu{અફસોસ}}, \emph{aphsos}.) If only? (\textbf{\gu{કાશ}}, \emph{kāś}.) To long for someone? (\textbf{\gu{ઝૂરવું}}, \emph{jhūrvũ}.) An opening? (\textbf{\gu{મોકો}}, \emph{moko}.) A goal? (\textbf{\gu{ધ્યેય}}, \emph{dhyey}.) A resolution? (\textbf{\gu{સંકલ્પ}}, \emph{sankalp}.) An achievement? (\textbf{\gu{સિદ્ધિ}}, \emph{siddhi}.) Planning? (\textbf{\gu{આયોજન}}, \emph{āyojan}.) To strive? (\textbf{\gu{મથવું}}, \emph{mathvũ}.) An attempt? (\textbf{\gu{પ્રયત્ન}}, \emph{prayatna}.) An effort? (\textbf{\gu{કોશિશ}}, \emph{kośiś}.) Enthusiasm? (\textbf{\gu{ઉત્સાહ}}, \emph{utsāh}.) To set about? (\textbf{\gu{માંડવું}}, \emph{mā̃ḍvũ}.) To battle on? (\textbf{\gu{ઝૂઝવું}}, \emph{jhūjhvũ}.) Successful? (\textbf{\gu{સફળ}}, \emph{saphaḷ}.) Advancement? (\textbf{\gu{ઉન્નતિ}}, \emph{unnati}.) Development? (\textbf{\gu{વિકાસ}}, \emph{vikās}.) To get one's way? (\textbf{\gu{ફાવવું}}, \emph{phāvvũ}.) To deliver? (\textbf{\gu{પહોંચાડવું}}, \emph{pahõcāḍvũ}.) An award? (\textbf{\gu{પુરસ્કાર}}, \emph{puraskār}.) Appreciation? (\textbf{\gu{કદર}}, \emph{kadar}.) Payment for work? (\textbf{\gu{મહેનતાણું}}, \emph{mahentāṇũ}.) To attain? (\textbf{\gu{પામવું}}, \emph{pāmvũ}.) To honour with an award? (\textbf{\gu{નવાજવું}}, \emph{navājvũ}.)
\end{tcolorbox}
